\section{Elektr sig`imi. Kondensatorlar}

\begin{enumerate}

\item O‘tkazgichning elektr sig‘imi deb ...
\begin{enumerate}[label=\Alph*)]
\item o‘tkazgich zaryadining potensialiga ko‘paytmasi bilan o‘lchanadigan kattalikka aytiladi.
\item o‘tkazgich zaryadiga to‘g‘ri proporsional va potensialiga teskari proporsional kattalikka aytiladi.
\item o‘tkazgich potensialining uning zaryadiga nisbatiga teng bo‘lgan kattalikka aytiladi.
\item o‘tkazgich zaryadining potensialiga nisbati bilan o‘lchanadigan kattalikka aytiladi.
\item TJY.
\end{enumerate}

\item Yakkalangan o‘tkazgichga $3\cdot10^{-9}$ C zaryad berilganda, uning potensiali $120$ V ga teng bo‘ldi. O‘tkazgichning elektr sig‘imini toping (F).
\begin{enumerate}[label=\Alph*)]
\item $2,5\cdot10^{-11}$
\item $4\cdot10^{-11}$
\item $3,6\cdot10^{-12}$
\item $3,6\cdot10^{-7}$
\item $5\cdot10^{-10}$
\end{enumerate}

\item Kondensatorning bir qoplamasining zaryadi $100$ nC, ikkinchisiniki esa $-100$ nC. Kondensatorning zaryadi necha nC?
\begin{enumerate}[label=\Alph*)]
\item 0
\item 100
\item 200
\item 1000
\item $10^4$
\end{enumerate}

\item Elektr sig‘imi $15$ nF bo‘lgan kondensatorning qoplamalari orasidagi potensiallar farqi $20$ V ga teng. Kondensatorning qoplamalaridan biridagi zaryad miqdorini toping (C).
\begin{enumerate}[label=\Alph*)]
\item 0,1
\item 0,3
\item 0,5
\item 3
\item 0,75
\end{enumerate}

\item $200$ V kuchlanishga ulangan, zaryadi $2\cdot10^{-4}$ C bo‘lgan kondensatorning sig‘imini toping ($\mu$F).
\begin{enumerate}[label=\Alph*)]
\item 0,1
\item 0,2
\item 0,4
\item 1
\item 2
\end{enumerate}

\item $300$ V o‘zgarmas kuchlanishga ulangan kondensator zaryadi $12\cdot10^{-4}$ C bo‘lsa, kondensatorning sig‘imi necha $\mu$F?
\begin{enumerate}[label=\Alph*)]
\item 0,1
\item 1
\item 2
\item 4
\item 7
\end{enumerate}

\item Sig‘imi $30$ pF bo‘lgan kondensator qoplamalariga berilgan zaryad miqdori $4,2\cdot10^{-8}$ C. Kondensator qanday kuchlanishgacha zaryadlangan (kV)?
\begin{enumerate}[label=\Alph*)]
\item 1,4
\item 34,2
\item 7
\item 2,8
\item 25,8
\end{enumerate}

\item Kondensatorning sig‘imi $4$ pF, qoplamalari orasidagi potensiallar ayirmasi $200$ V bo‘lsa, har bir qoplamada necha Kulon zaryad bor?
\begin{enumerate}[label=\Alph*)]
\item $6\cdot10^{-10}$
\item $2\cdot10^{-10}$
\item $2\cdot10^{-9}$
\item $8\cdot10^{-10}$
\item $4\cdot10^{-10}$
\end{enumerate}

\item Yassi kondensatorning bir plastinkasiga $+0,2$ C, ikkinchi plastinkasiga $-0,2$ C zaryad berilgan. Plastinkalar orasidagi kuchlanish $5$ kV. Kondensator sig‘imi qanday ($\mu$F)?
\begin{enumerate}[label=\Alph*)]
\item 80
\item 40
\item 25
\item 10
\item 0
\end{enumerate}

\item Kondensator qoplamalaridan biridagi potensial $300$ V, ikkinchisidagi esa $-300$ V va qoplamalardan biridagi zaryad $2,4$ $\mu$C bo‘lsa, uning sig‘imini toping (nF).
\begin{enumerate}[label=\Alph*)]
\item 8
\item 4
\item 5
\item 3
\item 10
\end{enumerate}

\item BXS dagi o‘lchamligi J/V$^2$ bo‘lgan fizik kattalik birligining nomi nima?
\begin{enumerate}[label=\Alph*)]
\item Kulon
\item Farad
\item Om
\item Amper
\item Nyuton
\end{enumerate}

\item Quyidagi birliklarning qaysilari elektr sig‘im birligi sifatida ishlatsa bo‘ladi:\\
1) F, \quad 2) C$^2$/J, 3) \quad 3) J/V$^2$, \quad 4) C/V.

\begin{enumerate}[label=\Alph*)]
\item 1, 4
\item 1, 2
\item hammasini
\item 1, 2, 3
\item 3, 4
\end{enumerate}

\item Yakkalangan shar elektr sig‘imi ifodasini ko‘rsating.
\begin{enumerate}[label=\Alph*)]
\item $C=\dfrac{\varphi_1-\varphi_2}{q}$
\item $C=\dfrac{\varepsilon\varepsilon_0S}{d}$
\item $C=2\pi\varepsilon\varepsilon_0R$
\item $C=4\pi\varepsilon\varepsilon_0R$
\item $C=4\pi\varepsilon R$
\end{enumerate}

\item Necha metr radiusli sharning elektr sig‘imi vakuumda $1$ F ga teng bo‘ladi?
\begin{enumerate}[label=\Alph*)]
\item $1\cdot10^7$
\item $1\cdot10^8$
\item $9\cdot10^{10}$
\item $1\cdot10^9$
\item $9\cdot10^9$
\end{enumerate}

\item O‘tkazgichga boshqa o‘tkazgich yaqinlashtirilganda, uning sig‘imi qanday o‘zgaradi?
\begin{enumerate}[label=\Alph*)]
\item ortadi
\item o‘zgarmaydi
\item kamayadi
\item avval ortadi, so‘ngra kamayadi
\item avval kamayadi, so‘ngra ortadi
\end{enumerate}

\item Zaryadlangan kondensator qoplamalari sim bilan tutashtirilib, zaryadsizlantirildi. Bunda qoplamalar massasi qanday o‘zgaradi?
\begin{enumerate}[label=\Alph*)]
\item har ikkala qoplamaning massasi ortadi.
\item musbat zaryadlangan qoplama massasi ortib, manfiy zaryadlangan qoplama massasi kamayadi.
\item har ikkala qoplamaning massasida o‘zgarish bo‘lmaydi.
\item manfiy zaryadlangan qoplama massasi ortib, musbat zaryadlangan qoplama massasi kamayadi.
\item har ikkala qoplamaning massasi kamayadi.
\end{enumerate}

\item Keltirilgan tenglamalardan qaysi biri yassi kondensator elektr sig‘imini ifodalaydi?
\begin{enumerate}[label=\Alph*)]
\item $C=C_1+C_2+\ldots+C_n$
\item $C=4\pi\varepsilon_0R$
\item $C=\dfrac{q}{U}$
\item $C=\dfrac{Q}{m(t_2-t_1)}$
\item $C=\dfrac{\varepsilon\varepsilon_0S}{d}$
\end{enumerate}

\item Yassi kondensator qoplamalarining yuzasi $40$ cm$^2$ dan, qoplamalari orasidagi masofa esa $5$ mm ga teng. Qoplamalar orasidagi muhitning dielektrik singdiruvchanligi $5$ ga teng bo‘lsa, uning sig‘imini toping (pF).
\begin{enumerate}[label=\Alph*)]
\item 12,46
\item 23,77
\item 15,1
\item 26,9
\item 35,4
\end{enumerate}

\item Sig‘imi $20$ pF bo‘lgan yassi kondensator har bir qoplamasining yuzasi $20$ cm$^2$ dan, qoplamalari orasidagi masofa esa $8,85$ mm ga teng. Qoplamalar orasidagi muhitning dielektrik singdiruvchanligini toping.
\begin{enumerate}[label=\Alph*)]
\item 8
\item 4
\item 5
\item 3
\item 10
\end{enumerate}

\item Sig‘imi $8,85$ pF bo‘lgan yassi kondensator har bir qoplamasining yuzasi $10$ cm$^2$ dan va qoplamalari orasidagi muhitning dielektrik singdiruvchanligi $3$ ga teng bo‘lsa, qoplamalari orasidagi masofani toping (mm).
\begin{enumerate}[label=\Alph*)]
\item 1
\item 4
\item 2
\item 3
\item 5
\end{enumerate}

\item Har bir qoplamasidagi zaryad miqdori $88,5$ $\mu$C bo‘lgan yassi kondensator qoplamalarining yuzasi $50$ cm$^2$ dan, qoplamalari orasidagi masofa $5$ mm va qoplamalari orasidagi muhitning dielektrik singdiruvchanligi $2,5$ ga teng bo‘lsa, qoplamalar orasidagi potensiallar farqini toping (MV).
\begin{enumerate}[label=\Alph*)]
\item 1
\item 4
\item 2
\item 3
\item 5
\end{enumerate}

\item Qoplamalari orasidagi kuchlanish $20$ kV bo‘lgan yassi kondensator qoplamalarining yuzasi $2$ cm$^2$ dan, qoplamalari orasidagi masofa $35,4$ mm va qoplamalari orasidagi muhitning dielektrik singdiruvchanligi $40$ ga teng bo‘lsa, har bir qoplamadagi zaryad miqdorini toping ($\mu$C).
\begin{enumerate}[label=\Alph*)]
\item 0,04
\item 0,03
\item 0,02
\item 0,15
\item 0,25
\end{enumerate}

\item Yassi kondensator qoplamalari diametri $0,2$ m dan bo‘lgan ikkita doiraviy plastinadan iborat. Qoplamalar orasidagi havo qatlamining qalinligi $3\cdot10^{-3}$ m, qoplamalar orasidagi kuchlanish $120$ V. Qoplamalardan har biridagi zaryad necha Kulon?
\begin{enumerate}[label=\Alph*)]
\item $1,55\cdot10^{-8}$
\item $1,22\cdot10^{-8}$
\item $0,82\cdot10^{-8}$
\item $1,01\cdot10^{-8}$
\item $1,06\cdot10^{-8}$
\end{enumerate}

\item Kondensator zaryadi 2 marta ortganda, uning sig‘imi qanday o‘zgaradi?
\begin{enumerate}[label=\Alph*)]
\item 4 marta kamayadi.
\item 2 marta oshadi.
\item 2 marta kamayadi.
\item 4 marta oshadi.
\item o‘zgarmaydi.
\end{enumerate}

\item Kondensator qoplamalaridagi zaryad 4 marta orttirilsa, uning sig‘imi qanday o‘zgaradi?
\begin{enumerate}[label=\Alph*)]
\item 4 marta kamayadi.
\item 2 marta ortadi.
\item 2 marta kamayadi.
\item 4 marta ortadi.
\item o‘zgarmaydi.
\end{enumerate}

\item Kondensator qoplamalari orasidagi kuchlanish 2 marta ortganda, uning sig‘imi qanday o‘zgaradi?
\begin{enumerate}[label=\Alph*)]
\item 2 marta kamayadi.
\item 2 marta ortadi.
\item o‘zgarmaydi.
\item 4 marta ortadi.
\item 4 marta kamayadi.
\end{enumerate}

\item Yassi kondensator qoplamasining yuzi 2 marta ortganda, uning sig‘imi qanday o‘zgaradi?
\begin{enumerate}[label=\Alph*)]
\item 2 marta kamayadi.
\item 4 marta ortadi.
\item o‘zgarmaydi.
\item 4 marta kamayadi.
\item 2 marta ortadi.
\end{enumerate}

\item Uchta yassi kondensator qoplamalarining yuzlari bir xil, ular orasidagi masofalar $d_1<d_2<d_3$ munosabatda bo‘lsa, sig‘imlarining munosabati qanday bo‘ladi?
\begin{enumerate}[label=\Alph*)]
\item $C_2<C_3<C_1$.
\item $C_1<C_2<C_3$.
\item $C_1=C_2=C_3$.
\item $C_1>C_3>C_2$.
\item $C_3<C_1<C_2$.
\end{enumerate}

\item Yassi havo kondensatori plastinalari orasidagi fazoga uchinchi plastina kiritilsa, sig‘imi qanday o‘zgaradi?
\begin{enumerate}[label=\Alph*)]
\item plastinaning dielektrik singdiruvchanligiga bog‘liq.
\item ortadi.
\item kamayadi.
\item o‘zgarmaydi.
\item TJY.
\end{enumerate}

\item Kondensator qoplamalari orasiga qalinligi qoplamalar orasidagi masofaning yarmiga teng bo‘lgan o‘tkazuvchi plastina kiritildi. Bunda kondensatorning sig‘imi qanday o‘zgaradi?
\begin{enumerate}[label=\Alph*)]
\item 4 marta ortadi.
\item 4 marta kamayadi.
\item o‘zgarmaydi.
\item 2 marta kamayadi.
\item 2 marta ortadi.
\end{enumerate}

\item Yassi kondensatorning har bir qoplamasi yuzi $S$, ular orasidagi masofa $d$. Qalinligi $l$ bo‘lgan metall plastina qoplamalarga tekkizilmagan holda kondensator ichiga kiritilgan. Kondensatorning sig‘imi qanday?
\begin{enumerate}[label=\Alph*)]
\item $\dfrac{\varepsilon_0 S}{2(d-l)}$
\item $\dfrac{\varepsilon_0 S}{d+l}$
\item $\dfrac{2\varepsilon_0 S}{d-l}$
\item $\dfrac{\varepsilon_0 S}{2(d+l)}$
\item $\dfrac{\varepsilon_0 S}{d-l}$
\end{enumerate}

\item Havo kondensatori qoplamalari orasidagi fazo hajmi $2$ m$^3$, uning sig‘imi $5$ nF. Agar hajmi $4$ m$^3$ gacha oshirilsa, kondensatorning sig‘imi qanday (nF) bo‘ladi?
\begin{enumerate}[label=\Alph*)]
\item 10
\item 15
\item 2,5
\item aniqlab bo‘lmaydi.
\item 5
\end{enumerate}

\item Yassi kondensator plastinalari orasidagi masofa $6$ marta, har bir plastinaning yuzi $3$ marta kamaytirilsa, kondensator sig‘imi qanday o‘zgaradi?
\begin{enumerate}[label=\Alph*)]
\item 2 marta ortadi.
\item 6 marta kamayadi.
\item 3 marta kamayadi.
\item 6 marta ortadi.
\item 3 marta ortadi.
\end{enumerate}

\item Yassi kondensator qoplamalari orasiga $\varepsilon=2$ bo‘lgan dielektrik bilan to‘ldirilsa, uning sig‘imi qanday o‘zgaradi?
\begin{enumerate}[label=\Alph*)]
\item o‘zgarmaydi.
\item 2 marta kamayadi.
\item 2 marta oshadi.
\item 4 marta kamayadi.
\item 4 marta oshadi.
\end{enumerate}

\item Uchta bir xil yassi kondensator qoplamalari orasiga turli dielektriklar joylashtirilgan: $\varepsilon_1<\varepsilon_2<\varepsilon_3$. Kondensatorlarning sig‘imlari qanday munosabatda bo‘ladi?
\begin{enumerate}[label=\Alph*)]
\item $C_3>C_2<C_1$.
\item $C_1>C_2>C_3$.
\item $C_1=C_2=C_3$.
\item $C_1<C_2<C_3$.
\item $C_1>C_2<C_3$.
\end{enumerate}

\item Yassi kondensator qoplamalari orasidagi masofa $4$ marta qisqartirilib, singdiruvchanligi $5$ ga teng bo‘lgan dielektrik o‘rniga, singdiruvchanligi $10$ ga teng bo‘lgan dielektrik joylashtirildi. Bunda kondensator sig‘imi qanday o‘zgaradi?
\begin{enumerate}[label=\Alph*)]
\item 4 marta kamayadi.
\item 8 marta ortadi.
\item 2 marta kamayadi.
\item 2 marta ortadi.
\item 8 marta kamayadi.
\end{enumerate}

\item O‘zgarmas tok manbaiga ulangan yassi kondensator plastinalari orasidagi masofa 2 marta oshirilsa, plastinalardagi zaryad miqdori qanday o‘zgaradi?
\begin{enumerate}[label=\Alph*)]
\item 4 marta kamayadi.
\item 2 marta oshadi.
\item o‘zgarmaydi.
\item 2 marta kamayadi.
\item 4 marta oshadi.
\end{enumerate}

\item Havo oraliqli yassi kondensator o‘zgarmas tok manbaidan zaryadlandi va undan uzildi. Plastinalar orasidagi masofa 2 marta kattalashtirilsa, kondensator zaryadi va kuchlanishi qanday o‘zgaradi?
\begin{enumerate}[label=\Alph*)]
\item zaryad va kuchlanish o‘zgarmaydi.
\item zaryad 2 marta ortadi, kuchlanish o‘zgarmaydi.
\item zaryad o‘zgarmaydi, kuchlanish 2 marta ortadi.
\item zaryad o‘zgarmaydi, kuchlanish 2 marta kamayadi.
\item zaryad 2 marta kamayadi, kuchlanish o‘zgarmaydi.
\end{enumerate}

\item Zaryadlangan kondensatorni manbadan uzib, plastinalari orasidagi masofa 7 marta kamaytirilsa, kuchlanishi qanday o‘zgaradi?
\begin{enumerate}[label=\Alph*)]
\item o‘zgarmaydi.
\item 7 marta ortadi.
\item 7 marta kamayadi.
\item 49 marta ortadi.
\item 49 marta kamayadi.
\end{enumerate}

\item Kondensator zaryadlanib, manbadan uzildi, so‘ngra u $\varepsilon=3$ bo‘lgan dielektrik bilan to‘ldirildi. Kondensatordagi kuchlanish qanday o‘zgaradi?
\begin{enumerate}[label=\Alph*)]
\item 9 marta ortadi.
\item 9 marta kamayadi.
\item o‘zgarmaydi.
\item 3 marta ortadi.
\item 3 marta kamayadi.
\end{enumerate}

\item Manbadan uzib qo‘yilgan zaryadlangan yassi kondensator plastinalari orasidagi masofa 2 marta oshirilsa, kondensatorning elektr maydon kuchlanganligi qanday o‘zgaradi?
\begin{enumerate}[label=\Alph*)]
\item o‘zgarmaydi.
\item 2 marta kamayadi.
\item 4 marta ortadi.
\item 4 marta kamayadi.
\item 2 marta ortadi.
\end{enumerate}

\item Agar zaryadlangan yassi havo kondensatorini manbadan uzib, qoplamalari orasidagi masofa 4 marta kamaytirilsa, undagi elektr maydon kuchlanganligi qanday o‘zgaradi?
\begin{enumerate}[label=\Alph*)]
\item o‘zgarmaydi.
\item 4 marta kamayadi.
\item 4 marta oshadi.
\item 16 marta kamayadi.
\item 16 marta oshadi.
\end{enumerate}

\item Yassi kondensator o‘zgarmas kuchlanish manbaiga ulandi. Qoplamalar orasidagi fazo dielektrik bilan to‘ldirilsa, kondensator ichidagi elektr maydon kuchlanganligi qanday o‘zgaradi?
\begin{enumerate}[label=\Alph*)]
\item kamayadi.
\item ortadi.
\item o‘zgarmaydi.
\item ortishi ham mumkin, kamayishi ham.
\item javob dielektrikning singdiruvchanligiga bog‘liq.
\end{enumerate}

\item Zaryadlangan kondensator qoplamalari bir-biriga $F$ kuch bilan tortiladilar. Agar qoplamalar orasiga dielektrik singdiruvchanligi $\varepsilon=2$ bo‘lgan yassi shisha plastina kiritilsa, bu kuch qanday o‘zgaradi? Plastina qalinligi qoplamalar orasidagi masofadan kichik bo‘lsin.
\begin{enumerate}[label=\Alph*)]
\item 2 marta kamayadi.
\item 2 marta ortadi.
\item 4 marta kamayadi.
\item 4 marta ortadi.
\item o‘zgarmaydi.
\end{enumerate}

\item Yassi havo kondensatori qoplamalari orasidagi masofa $8,85$ mm, zaryadining sirt zichligi $10^{-6}$ C/m$^2$ bo‘lsa, uning kuchlanishini necha Volt bo‘ladi?
\begin{enumerate}[label=\Alph*)]
\item 100
\item 500
\item 885
\item 1000
\item 2000
\end{enumerate}

\item Yassi havo kondensatori $9$ kV kuchlanishga ega bo‘lgan tok manbaiga ulangan. Kondensator plastinkalari orasidagi maydon kuchlanganligi $3$ MV/m bo‘lganda, havoda ionizatsiya tufayli teshilish yuz bersa, kondensator ishdan chiqmasligi uchun uning qoplamalarini qanday minimal masofaga yaqinlashtirish mumkin (mm)?
\begin{enumerate}[label=\Alph*)]
\item 1,5
\item 3
\item 15
\item 30
\item 60
\end{enumerate}


\end{enumerate}
\begin{table}[h!]
\centering
\begin{tabular}{|c|c|c|c|c|c|c|c|c|c|}
\hline
1d & 2a & 3b & 4b & 5d & 6d & 7a & 8d & 9b & 10b \\ \hline
11b & 12c & 13d & 14e & 15a & 16b & 17e & 18e & 19e & 20d \\ \hline
21b & 22a & 23e & 24e & 25e & 26c & 27e & 28d & 29b & 30e \\ \hline
31e & 32d & 33a & 34c & 35d & 36b & 37d & 38c & 39c & 40e \\ \hline
41a & 42a & 43c & 44e & 45d & 46b &  &  &  &  \\ \hline
\end{tabular}
\end{table}